\documentclass[10pt,a4paper]{amsart}
\usepackage{amsmath,amssymb,amsthm,graphicx}
\usepackage{verbatim}
\newenvironment{reference}{\comment}{\endcomment}
%\newenvironment{reference}{}{}
\newenvironment{slogan}{\comment}{\endcomment}
\newenvironment{history}{\comment}{\endcomment}

% For commutative diagrams we use Xy-pic
\usepackage[all]{xy}

% We use 2cell for 2-commutative diagrams.
\xyoption{2cell}
% Theorem environments.
%
\theoremstyle{plain}
\newtheorem{theorem}[subsection]{Theorem}
\newtheorem{proposition}[subsection]{Proposition}
\newtheorem{lemma}[subsection]{Lemma}

\theoremstyle{definition}
\newtheorem{definition}[subsection]{Definition}
\newtheorem{example}[subsection]{Example}
\newtheorem{exercise}[subsection]{Exercise}
\newtheorem{situation}[subsection]{Situation}

\theoremstyle{remark}
\newtheorem{remark}[subsection]{Remark}
\newtheorem{remarks}[subsection]{Remarks}

\numberwithin{equation}{subsection}

% Macros
%
\def\lim{\mathop{\mathrm{lim}}\nolimits}
\def\colim{\mathop{\mathrm{colim}}\nolimits}
\def\Spec{\mathop{\mathrm{Spec}}}
\def\Hom{\mathop{\mathrm{Hom}}\nolimits}
\def\Ext{\mathop{\mathrm{Ext}}\nolimits}
\def\SheafHom{\mathop{\mathcal{H}\!\mathit{om}}\nolimits}
\def\SheafExt{\mathop{\mathcal{E}\!\mathit{xt}}\nolimits}
\def\Sch{\mathit{Sch}}
\def\Mor{\mathop{\mathrm{Mor}}\nolimits}
\def\Ob{\mathop{\mathrm{Ob}}\nolimits}
\def\Sh{\mathop{\mathit{Sh}}\nolimits}
\def\NL{\mathop{N\!L}\nolimits}
\def\CH{\mathop{\mathrm{CH}}\nolimits}
\def\proetale{{pro\text{-}\acute{e}tale}}
\def\etale{{\acute{e}tale}}
\def\QCoh{\mathit{QCoh}}
\def\Ker{\mathop{\mathrm{Ker}}}
\def\Im{\mathop{\mathrm{Im}}}
\def\Coker{\mathop{\mathrm{Coker}}}
\def\Coim{\mathop{\mathrm{Coim}}}

% Boxtimes
%
\DeclareMathSymbol{\boxtimes}{\mathbin}{AMSa}{"02}

%
% Macros for moduli stacks/spaces
%
\def\QCohstack{\mathcal{QC}\!\mathit{oh}}
\def\Cohstack{\mathcal{C}\!\mathit{oh}}
\def\Spacesstack{\mathcal{S}\!\mathit{paces}}
\def\Quotfunctor{\mathrm{Quot}}
\def\Hilbfunctor{\mathrm{Hilb}}
\def\Curvesstack{\mathcal{C}\!\mathit{urves}}
\def\Polarizedstack{\mathcal{P}\!\mathit{olarized}}
\def\Complexesstack{\mathcal{C}\!\mathit{omplexes}}
% \Pic is the operator that assigns to X its picard group, usage \Pic(X)
% \Picardstack_{X/B} denotes the Picard stack of X over B
% \Picardfunctor_{X/B} denotes the Picard functor of X over B
\def\Pic{\mathop{\mathrm{Pic}}\nolimits}
\def\Picardstack{\mathcal{P}\!\mathit{ic}}
\def\Picardfunctor{\mathrm{Pic}}
\def\Deformationcategory{\mathcal{D}\!\mathit{ef}}
\title{An inverse system of quasi-compact topological spaces whose limit is not quasi-compact}
\author{The Stacks Project Authors}
\date{stacks-project:examples.tex:local-snapshot-20261012; GNU FDL 1.2 or later}
\begin{document}
\maketitle
\noindent\textit{Editorial note.} All delivered mathematical text is the author's own, copied from the pinned Stacks Project chapter file examples.tex; no statement or proof step was rewritten, reordered or removed.\par\medskip
\noindent\textit{Editorial note.} The counted proof is the author's complete argument: The counted proof adds no mathematics of its own; the singleton covering having no finite subcover is the context's closing observation (333-335).\par\medskip
\noindent\textit{Editorial note.} The label lemma-lim-not-quasi-compact names the counted statement, and the proof environment immediately following it in the source is delivered verbatim as its proof.\par\medskip
\noindent\textit{Editorial note.} The counted statement is carried by the author's own lemma environment in the chapter source.\par\medskip
\noindent\textit{Allowed standard prerequisites (editorial).}\begin{itemize}
\item Context span 296-335: for each n the topology T\_n on N with basis \{1\}, ..., \{n\}, I\_n where I\_n = naturals > n; each X\_n = (N, T\_n) is quasi-compact and the identity maps f\_\{n,m\}: X\_n -> X\_m are continuous with the compatibilities f\_\{p,n\} compose f\_\{n,m\} = f\_\{p,m\}, giving a directed inverse system of quasi-compact spaces
\item Context continues: X = (N, discrete) with f\_n: X -> X\_n continuous; the universal property check (any compatible system of maps g\_n: (Y, S) -> (N, T\_n) is a common set map g: Y -> N with all g\textasciicircum{}\{-1\}(\{n\}) open in Y, hence continuous for S and T) shows X is the inverse limit; finally the infinite open covering of X by singleton sets has no inverse limit, so X is not quasi-compact
\end{itemize}\medskip
Let $\mathbf{N}$ denote the set of natural numbers.
For every integer $n$, let $I_n$ denote the set of all natural numbers $> n$.
Define $T_n$ to be the unique topology on $\mathbf{N}$ with basis
$\{1\}, \ldots , \{n\}, I_n$.  Denote by $X_n$ the topological space
$(\mathbf{N}, T_n)$.  For each $m < n$, the identity map,
$$
f_{n, m} : X_n \longrightarrow X_m
$$
is continuous.  Obviously for $m < n < p$, the composition
$f_{p, n} \circ f_{n, m}$ equals $f_{p, m}$.  So $((X_n), (f_{n,m}))$
is a directed inverse system of quasi-compact topological spaces.

\medskip\noindent
Let $T$ be the discrete topology on $\mathbf{N}$, and let $X$ be
$(\mathbf{N}, T)$. Then for every integer $n$, the identity map,
$$
f_n : X \longrightarrow X_n
$$
is continuous. We claim that this is the inverse limit of the directed
system above. Let $(Y, S)$ be any topological space. For every integer $n$,
let
$$
g_n : (Y, S) \longrightarrow (\mathbf{N}, T_n)
$$
be a continuous map. Assume that for every $m < n$ we have
$f_{n,m} \circ g_n = g_m$, i.e., the system $(g_n)$ is compatible
with the directed system above. In particular, all of the set maps
$g_n$ are equal to a common set map
$$
g : Y \longrightarrow \mathbf{N}.
$$
Moreover, for every integer $n$, since $\{n\}$ is open in $X_n$,
also $g^{-1}(\{n\}) = g_n^{-1}(\{n\})$ is open in $Y$.
Therefore the set map $g$ is continuous for the topology $S$ on $Y$
and the topology $T$ on $\mathbf{N}$. Thus $(X, (f_n))$ is the inverse
limit of the directed system above.

\medskip\noindent
However, clearly $X$ is not quasi-compact, since the infinite open
covering by singleton sets has no inverse limit.

\begin{lemma}
\label{lemma-lim-not-quasi-compact}
There exists an inverse system of quasi-compact topological spaces
over $\mathbf{N}$ whose limit is not quasi-compact.
\end{lemma}

\begin{proof}
See discussion above.
\end{proof}

\end{document}
